\documentclass[12pt]{article}

\usepackage[margin=1in]{geometry}
\usepackage{enumitem}
\usepackage{amsmath}
\usepackage{amssymb}
\usepackage{array}
\usepackage{titlesec}

\title{Lecture 4 Notes:\\ Aristotle's \textit{Prior Analytics}, Book I}
\author{}
\date{}

\begin{document}

\maketitle

\section*{Central Theme}

Aristotle's \textit{Prior Analytics}, Book I constructs the theory of the
syllogism. It explains what a syllogism is, how valid inference is structured,
how the figures work, how arguments are discovered, and how arguments can be
analyzed back into their logical form.

\[
\boxed{
\textit{Prior Analytics}, \text{ Book I}
=
\text{Aristotle's construction of the syllogism}
}
\]

\section*{1. From Terms to Propositions to Syllogisms}

The \textit{Prior Analytics} continues the movement of Aristotle's logical
works.

\[
\textit{Categories} \rightarrow \text{terms}
\]

\[
\textit{On Interpretation} \rightarrow \text{propositions}
\]

\[
\textit{Prior Analytics} \rightarrow \text{syllogisms}
\]

In the \textit{Categories}, Aristotle classified what can be said. In
\textit{On Interpretation}, he explained how sayings become true or false. In
the \textit{Prior Analytics}, he explains how propositions combine so that a
conclusion follows necessarily.

\subsection*{Teaching Point}

The syllogism is Aristotle's basic model of valid inference.

\[
\boxed{
\text{terms combine into propositions; propositions combine into syllogisms.}
}
\]

\section*{2. The Subject of the Inquiry}

Aristotle begins by saying that the subject of the inquiry is demonstration and
that the faculty which carries it out is demonstrative science.

But before demonstration can be studied, Aristotle must first explain
syllogism. Demonstration is a special kind of syllogism, but not every
syllogism is a demonstration.

\[
\text{demonstration} \subset \text{syllogism}
\]

\subsection*{Teaching Point}

Aristotle studies syllogism before demonstration because syllogism is more
general.

\[
\boxed{
\text{Every demonstration is a syllogism, but not every syllogism is a demonstration.}
}
\]

\section*{3. Premiss}

A premiss is a sentence affirming or denying one thing of another.

\[
\boxed{
\text{premiss} = \text{a sentence affirming or denying one thing of another}
}
\]

Premisses may be universal, particular, or indefinite.

\subsection*{Universal Premiss}

A universal premiss says that something belongs to all or none of something.

\[
\text{All humans are animals.}
\]

\[
\text{No stones are animals.}
\]

\subsection*{Particular Premiss}

A particular premiss says that something belongs to some, does not belong to
some, or does not belong to all.

\[
\text{Some animals are white.}
\]

\[
\text{Some animals are not human.}
\]

\[
\text{Not all animals are human.}
\]

\subsection*{Indefinite Premiss}

An indefinite premiss affirms or denies without marking whether the statement
is universal or particular.

\[
\text{Pleasure is not good.}
\]

\[
\text{Contraries are subjects of the same science.}
\]

\subsection*{Whiteboard}

\[
\begin{array}{ll}
\textbf{Universal} & \text{all / none} \\
\textbf{Particular} & \text{some / not some / not all} \\
\textbf{Indefinite} & \text{no explicit mark of quantity}
\end{array}
\]

\section*{4. Demonstrative and Dialectical Premisses}

Aristotle distinguishes demonstrative premisses from dialectical premisses.

A demonstrative premiss is true and obtained through the first principles of
its science.

A dialectical premiss depends on a choice between contradictories and proceeds
from what is apparent or generally admitted.

\[
\text{demonstrative premiss} = \text{true and scientific}
\]

\[
\text{dialectical premiss} = \text{accepted through questioning or common opinion}
\]

\subsection*{Teaching Point}

Demonstration and dialectic differ in the status of their premisses, but both
argue syllogistically.

\[
\boxed{
\text{The same syllogistic structure can serve different kinds of inquiry.}
}
\]

\section*{5. Term}

A term is that into which a premiss is resolved: the predicate and the subject
of which it is predicated.

\[
\boxed{
\text{term} = \text{predicate or subject within a premiss}
}
\]

For example:

\[
\text{Every human is animal.}
\]

The terms are:

\[
\text{human}, \quad \text{animal}
\]

\subsection*{Teaching Point}

A syllogism is built from terms arranged in propositions.

\[
\text{terms} \rightarrow \text{premisses} \rightarrow \text{conclusion}
\]

\section*{6. Syllogism}

Aristotle's definition of syllogism is one of the central definitions of the
\textit{Prior Analytics}.

\[
\boxed{
\text{A syllogism is discourse in which, certain things being stated,}
}
\]

\[
\boxed{
\text{something other than what is stated follows of necessity.}
}
\]

There are two important parts of this definition.

\subsection*{Something Other}

The conclusion must not simply repeat a premiss. It must be something other
than what has already been stated.

\subsection*{Follows of Necessity}

The conclusion follows because of the structure of the premisses. It is not
merely a true statement added after other true statements.

\subsection*{Teaching Point}

A syllogism is not just a sequence of propositions. It is a necessary inference.

\[
\boxed{
\text{Syllogism} = \text{structured necessity in reasoning}
}
\]

\section*{7. Perfect and Imperfect Syllogisms}

Aristotle distinguishes perfect and imperfect syllogisms.

\subsection*{Perfect Syllogism}

A perfect syllogism needs nothing beyond the stated premisses to make the
necessity of the conclusion clear.

\[
\boxed{
\text{perfect syllogism} = \text{complete from the premisses as stated}
}
\]

\subsection*{Imperfect Syllogism}

An imperfect syllogism is valid, but it needs one or more additional
propositions to make the conclusion clear.

\[
\boxed{
\text{imperfect syllogism} = \text{valid, but needing completion}
}
\]

\subsection*{Teaching Point}

``Imperfect'' does not mean invalid. It means that the syllogism must be
completed or reduced to a clearer form.

\[
\boxed{
\text{The first figure will become the standard of completion.}
}
\]

\section*{8. Predicated of All and Predicated of None}

Aristotle defines what it means for one term to be predicated of all of
another.

\[
\text{A is predicated of all B}
\]

means:

\[
\text{there is no instance of B of which A cannot be asserted.}
\]

Similarly:

\[
\text{A is predicated of no B}
\]

means:

\[
\text{there is no instance of B of which A can be asserted.}
\]

\subsection*{Whiteboard}

\[
A \text{ belongs to all } B
\quad \Longleftrightarrow \quad
\text{Every } B \text{ is } A
\]

\[
A \text{ belongs to no } B
\quad \Longleftrightarrow \quad
\text{No } B \text{ is } A
\]

\section*{9. Conversion of Pure Propositions}

Before Aristotle gives the figures, he explains conversion.

Conversion is the rearrangement of subject and predicate while preserving a
valid consequence.

\subsection*{Universal Negative Converts Universally}

\[
\text{No } A \text{ is } B
\quad \Longleftrightarrow \quad
\text{No } B \text{ is } A
\]

Example:

\[
\text{No pleasure is good.}
\]

Therefore:

\[
\text{No good is pleasure.}
\]

\subsection*{Universal Affirmative Converts Particularly}

\[
\text{All } A \text{ is } B
\quad \Longrightarrow \quad
\text{Some } B \text{ is } A
\]

Example:

\[
\text{Every pleasure is good.}
\]

Therefore:

\[
\text{Some good is pleasure.}
\]

\subsection*{Particular Affirmative Converts Particularly}

\[
\text{Some } A \text{ is } B
\quad \Longleftrightarrow \quad
\text{Some } B \text{ is } A
\]

Example:

\[
\text{Some pleasure is good.}
\]

Therefore:

\[
\text{Some good is pleasure.}
\]

\subsection*{Particular Negative Does Not Convert}

\[
\text{Some } A \text{ is not } B
\quad \not\Longrightarrow \quad
\text{Some } B \text{ is not } A
\]

Example:

\[
\text{Some animals are not humans.}
\]

It does not follow that:

\[
\text{Some humans are not animals.}
\]

\subsection*{Whiteboard}

\[
\begin{array}{lll}
E & \text{No } A \text{ is } B & \text{converts universally} \\
A & \text{All } A \text{ is } B & \text{converts particularly} \\
I & \text{Some } A \text{ is } B & \text{converts particularly} \\
O & \text{Some } A \text{ is not } B & \text{does not convert}
\end{array}
\]

\subsection*{Teaching Point}

Conversion is crucial because Aristotle uses it to reduce imperfect syllogisms
to the first figure.

\[
\boxed{
\text{Conversion is the tool of reduction.}
}
\]

\section*{10. The Middle Term}

The middle term is the term that connects the two extremes.

\[
\boxed{
\text{middle term} = \text{the connector between subject and predicate}
}
\]

Example:

\[
\text{All humans are animals.}
\]

\[
\text{All Greeks are humans.}
\]

\[
\therefore \text{All Greeks are animals.}
\]

Here the middle term is:

\[
\text{humans}
\]

It connects:

\[
\text{Greeks}
\quad \text{and} \quad
\text{animals}
\]

\subsection*{Teaching Point}

To understand or construct a syllogism, find the middle term.

\[
\boxed{
\text{No middle term, no syllogism.}
}
\]

\section*{11. The First Figure}

The first figure is the clearest and most important syllogistic figure.

In the first figure, the middle term stands between the extremes. The major
term is predicated of the middle, and the middle is predicated of the minor.

\[
A \text{ belongs to all } B
\]

\[
B \text{ belongs to all } C
\]

\[
\therefore A \text{ belongs to all } C
\]

\subsection*{Example}

\[
\text{All humans are animals.}
\]

\[
\text{All Greeks are humans.}
\]

\[
\therefore \text{All Greeks are animals.}
\]

\subsection*{Diagram}

\[
C \subseteq B \subseteq A
\]

\[
\therefore C \subseteq A
\]

\subsection*{Teaching Point}

The first figure is perfect because the necessity of the conclusion is evident
from the premisses themselves.

\[
\boxed{
\text{The first figure is the model of syllogistic clarity.}
}
\]

\section*{12. First Figure Conclusions}

The first figure can prove all four basic kinds of conclusions:

\[
\begin{array}{ll}
\textbf{Universal affirmative} & \text{All } C \text{ is } A \\
\textbf{Universal negative} & \text{No } C \text{ is } A \\
\textbf{Particular affirmative} & \text{Some } C \text{ is } A \\
\textbf{Particular negative} & \text{Some } C \text{ is not } A
\end{array}
\]

\subsection*{Teaching Point}

Because the first figure can prove every type of conclusion, it has a special
place in Aristotle's logical system.

\[
\boxed{
\text{All conclusions are provable through the first figure.}
}
\]

\section*{13. The Second Figure}

In the second figure, the middle term is predicated of both extremes.

\[
M \text{ belongs to no } N
\]

\[
M \text{ belongs to all } O
\]

\[
\therefore N \text{ belongs to no } O
\]

\subsection*{Structure}

\[
M \rightarrow N
\]

\[
M \rightarrow O
\]

The middle term stands outside the extremes and is predicate in both premisses.

\subsection*{Example}

\[
\text{No stone is alive.}
\]

\[
\text{Every plant is alive.}
\]

\[
\therefore \text{No plant is a stone.}
\]

\subsection*{Teaching Point}

The second figure produces negative conclusions and is imperfect. It must be
completed by conversion or reduction.

\[
\boxed{
\text{Second figure syllogisms depend on reduction to the first figure.}
}
\]

\section*{14. The Third Figure}

In the third figure, the middle term is subject of both predications. Both
extremes are predicated of the middle.

\[
P \text{ belongs to all } S
\]

\[
R \text{ belongs to all } S
\]

\[
\therefore P \text{ belongs to some } R
\]

\subsection*{Example}

\[
\text{All horses are animals.}
\]

\[
\text{All horses are living things.}
\]

\[
\therefore \text{Some living things are animals.}
\]

\subsection*{Teaching Point}

The third figure produces particular conclusions. It cannot produce a universal
conclusion.

\[
\boxed{
\text{The third figure shows that shared predication yields particular connection.}
}
\]

\section*{15. The Three Figures Compared}

The figure of a syllogism is determined by the position of the middle term.

\[
\boxed{
\text{The figure is determined by the position of the middle term.}
}
\]

\[
\begin{array}{c|c|c}
\textbf{Figure} & \textbf{Position of Middle Term} & \textbf{Typical Role} \\
\hline
\text{First} & \text{subject once, predicate once} & \text{direct proof} \\
\text{Second} & \text{predicate in both premisses} & \text{negative conclusions} \\
\text{Third} & \text{subject in both premisses} & \text{particular conclusions}
\end{array}
\]

\section*{16. Reduction to the First Figure}

A major claim of Book I is that imperfect syllogisms are completed through the
first figure.

\[
\text{Second figure} \rightarrow \text{First figure}
\]

\[
\text{Third figure} \rightarrow \text{First figure}
\]

This may happen by:

\begin{itemize}[itemsep=0.25em]
    \item conversion,
    \item exposition,
    \item reductio ad impossibile.
\end{itemize}

\subsection*{Teaching Point}

The first figure is not merely one figure among others. It is the standard to
which the others are reduced.

\[
\boxed{
\text{The first figure perfects the imperfect figures.}
}
\]

\section*{17. Reductio ad Impossibile}

Aristotle also uses reduction to impossibility.

In a reductio argument, one assumes the contradictory of what one wants to
prove. If an impossibility follows, the original conclusion is established.

\[
\text{Assume not-}P
\]

\[
\text{Derive impossibility}
\]

\[
\therefore P
\]

\subsection*{Teaching Point}

Even arguments by impossibility depend on syllogistic structure. They too can
be analyzed through the figures.

\[
\boxed{
\text{Reductio is indirect proof, but it still relies on syllogistic necessity.}
}
\]

\section*{18. Modal Syllogisms}

A large part of Book I concerns modal syllogisms.

Aristotle distinguishes premisses according to whether something:

\[
\text{belongs}
\]

\[
\text{belongs necessarily}
\]

\[
\text{may possibly belong}
\]

\subsection*{Three Modes of Attribution}

\[
\begin{array}{ll}
\textbf{Assertoric} & \text{is / belongs} \\
\textbf{Necessary} & \text{must be / necessarily belongs} \\
\textbf{Problematic} & \text{may be / possibly belongs}
\end{array}
\]

\subsection*{Teaching Point}

The modality of the premisses affects the modality of the conclusion.

\[
\boxed{
\text{From what is necessary, we may get a necessary conclusion.}
}
\]

\[
\boxed{
\text{From what is possible, we must be careful about what kind of possibility follows.}
}
\]

\section*{19. Necessary Premisses}

Aristotle says there is little difference between syllogisms from necessary
premisses and syllogisms from simple assertoric premisses, so long as the terms
are arranged in the same way.

The main difference is that necessity is added to the conclusion when the
appropriate premiss is necessary.

\subsection*{Example}

\[
\text{All humans are necessarily animals.}
\]

\[
\text{All Greeks are humans.}
\]

\[
\therefore \text{All Greeks are necessarily animals.}
\]

\subsection*{Important Qualification}

If only the minor premiss is necessary, the conclusion need not be necessary.

Example pattern:

\[
A \text{ belongs to } B
\]

\[
B \text{ necessarily belongs to } C
\]

It does not always follow that:

\[
A \text{ necessarily belongs to } C
\]

\subsection*{Teaching Point}

Necessity usually follows the major term's relation to the middle.

\[
\boxed{
\text{A necessary minor premiss alone does not guarantee a necessary conclusion.}
}
\]

\section*{20. Possible Premisses}

Aristotle defines the possible as that which is not necessary but, being
assumed, results in nothing impossible.

\[
\boxed{
\text{possible} = \text{not necessary, but not impossible}
}
\]

The possible is said in more than one way.

\subsection*{Natural Possibility}

Some things happen generally or naturally, though not necessarily.

Example:

\[
\text{A human grows grey.}
\]

\subsection*{Indefinite Possibility}

Some things can happen either way and incline no more to one side than to the
other.

Example:

\[
\text{An animal walks.}
\]

\[
\text{An earthquake occurs while it walks.}
\]

\subsection*{Teaching Point}

Aristotle is not merely asking whether something can happen. He is asking what
kind of possibility is involved.

\[
\boxed{
\text{Possibility is not one simple notion.}
}
\]

\section*{21. Why Modal Logic Is Difficult}

With modal syllogisms, the placement of necessity and possibility matters.

\[
\text{It may be}
\neq
\text{It must be}
\]

\[
\text{It may not be}
\neq
\text{It cannot be}
\]

\[
\text{Not necessarily belongs}
\neq
\text{Necessarily does not belong}
\]

\subsection*{Teaching Point}

Modal syllogisms require special care because different modes of predication
do not always convert or combine in the same way.

\[
\boxed{
\text{Changing the modality changes the inference.}
}
\]

\section*{22. Discovery of Arguments}

After classifying syllogisms, Aristotle turns to the discovery of arguments.

The question is:

\begin{quote}
How do we find the premisses needed for a syllogism?
\end{quote}

The answer is:

\[
\boxed{
\text{Find the middle term.}
}
\]

\subsection*{What to Gather}

For the subject under investigation, Aristotle tells us to collect:

\begin{itemize}[itemsep=0.25em]
    \item definitions,
    \item properties,
    \item accidents,
    \item what follows from the subject,
    \item what the subject follows from,
    \item what cannot belong to the subject.
\end{itemize}

\subsection*{Whiteboard}

\[
\text{subject}
\]

\[
\downarrow
\]

\[
\text{definitions, properties, accidents}
\]

\[
\downarrow
\]

\[
\text{what follows / what it follows / what cannot belong}
\]

\[
\downarrow
\]

\[
\text{middle term}
\]

\subsection*{Teaching Point}

Aristotle's logic is not only a theory of valid forms. It is also a method for
finding arguments.

\[
\boxed{
\text{To discover an argument, discover the middle.}
}
\]

\section*{23. The Middle Term as Cause of Inference}

In a syllogism, the middle term explains why the extremes are connected.

\[
\text{Major term} \longleftrightarrow \text{Middle term}
\longleftrightarrow \text{Minor term}
\]

Example:

\[
\text{All humans are animals.}
\]

\[
\text{All Greeks are humans.}
\]

\[
\therefore \text{All Greeks are animals.}
\]

The middle term is:

\[
\text{human}
\]

It explains why Greeks fall under animal.

\subsection*{Teaching Point}

The middle term is not merely a repeated word. It is the logical bridge.

\[
\boxed{
\text{The middle term makes the conclusion necessary.}
}
\]

\section*{24. Analysis of Arguments}

Aristotle then turns to the analysis of arguments into figures and moods.

To analyze an argument, we must:

\begin{enumerate}[itemsep=0.25em]
    \item Find the conclusion.
    \item Find the two premisses.
    \item Remove unnecessary assumptions.
    \item Supply missing necessary assumptions.
    \item Identify the terms.
    \item Find the middle term.
    \item Determine the figure.
    \item Check the quantity and quality of the premisses.
\end{enumerate}

\subsection*{Whiteboard}

\[
\boxed{
\text{To analyze an argument, find the middle term and the figure.}
}
\]

\section*{25. Necessary but Not Syllogistic}

Aristotle warns that not everything necessary is a syllogism.

Sometimes something follows necessarily from what has been assumed, but the
premisses are not arranged in proper syllogistic form.

\[
\boxed{
\text{Every syllogism is necessary, but not everything necessary is a syllogism.}
}
\]

\subsection*{Teaching Point}

A valid syllogism requires not only necessity, but a determinate arrangement of
terms and premisses.

\section*{26. How to Identify the Figure}

The figure is determined by the role of the middle term.

\[
\begin{array}{c|c}
\textbf{Figure} & \textbf{Middle Term} \\
\hline
\text{First} & \text{predicate in one premiss, subject in the other} \\
\text{Second} & \text{predicate in both premisses} \\
\text{Third} & \text{subject in both premisses}
\end{array}
\]

\subsection*{Teaching Point}

The same conclusion may be approached differently depending on how the middle
term is positioned.

\[
\boxed{
\text{The middle term determines the figure.}
}
\]

\section*{27. Division as a Weak Syllogism}

Aristotle criticizes the method of division.

He says division is like a weak syllogism because it tends to assume what it
ought to prove.

\subsection*{Example}

Suppose one divides animals into mortal and immortal, then places man under
animal. This may show that man is either mortal or immortal, but it does not
prove that man is mortal unless that is assumed.

\[
\text{Animal} \rightarrow \text{mortal or immortal}
\]

\[
\text{Man is animal}
\]

\[
\therefore \text{Man is mortal or immortal}
\]

But this does not yet prove:

\[
\text{Man is mortal}
\]

\subsection*{Teaching Point}

Division classifies, but syllogism proves.

\[
\boxed{
\text{Classification is not yet demonstration.}
}
\]

\section*{28. Ordinary Language and Logical Form}

In the later chapters of Book I, Aristotle gives many cautions about translating
ordinary language into syllogistic form.

He warns that we must:

\begin{itemize}[itemsep=0.25em]
    \item set out terms correctly,
    \item use equivalent expressions when needed,
    \item avoid being misled by grammatical form,
    \item understand different kinds of predication,
    \item distinguish simple from qualified predication.
\end{itemize}

\subsection*{Teaching Point}

Logic requires interpretation. We cannot always move directly from ordinary
grammar to syllogistic form.

\[
\boxed{
\text{Grammatical form and logical form are not always identical.}
}
\]

\section*{29. ``Is Not A'' and ``Is Not-A''}

Book I ends with an important distinction.

\[
\text{not to be white}
\neq
\text{to be not-white}
\]

The negation of:

\[
\text{to be white}
\]

is:

\[
\text{not to be white}
\]

not:

\[
\text{to be not-white}
\]

\subsection*{Example}

\[
\text{This is not a white log.}
\]

does not mean the same thing as:

\[
\text{This is a not-white log.}
\]

The first could be true of something that is not a log at all. The second
requires that the thing be a log.

\subsection*{Teaching Point}

Negation and indefinite predication are different.

\[
\boxed{
\text{Negation} \neq \text{indefinite predicate}
}
\]

\section*{30. The Architecture of Book I}

Book I has three broad movements.

\[
\begin{array}{ll}
\textbf{I. Structure of the syllogism} &
\text{definitions, conversion, figures, modal syllogisms} \\
\textbf{II. Discovery of arguments} &
\text{finding middle terms and premisses} \\
\textbf{III. Analysis of arguments} &
\text{reducing arguments to figures and moods}
\end{array}
\]

\subsection*{Teaching Point}

Book I is not merely a list of valid forms. It is a complete logical technology.

\[
\boxed{
\text{Aristotle teaches how to build, find, and analyze arguments.}
}
\]

\section*{31. Summary of Main Ideas}

\begin{enumerate}[itemsep=0.5em]
    \item A syllogism is a discourse in which something other than the stated
    premisses follows of necessity.

    \item A premiss affirms or denies one thing of another.

    \item A term is what a premiss is resolved into: subject and predicate.

    \item The middle term connects the two extremes.

    \item The first figure is the clearest and most complete figure.

    \item The second and third figures are imperfect and are completed through
    the first.

    \item Modal syllogisms require attention to necessity and possibility.

    \item The discovery of arguments depends on finding the right middle term.

    \item The analysis of arguments depends on identifying premisses, terms,
    middle term, figure, and mood.

    \item Ordinary language must be carefully interpreted before it can be
    placed into syllogistic form.
\end{enumerate}

\section*{Final Framing}

\[
\boxed{
\textit{Prior Analytics}, \text{ Book I, gives Aristotle's machinery of inference.}
}
\]

In the \textit{Categories}, Aristotle classified terms. In
\textit{On Interpretation}, he explained propositions. In the
\textit{Prior Analytics}, he shows when propositions force a conclusion.

\[
\boxed{
\text{A syllogism is the moment when propositions become necessary reasoning.}
}
\]

\end{document}